\documentclass{article}%
\usepackage[T1]{fontenc}%
\usepackage[utf8]{inputenc}%
\usepackage{lmodern}%
\usepackage{textcomp}%
\usepackage{lastpage}%
\usepackage{geometry}%
\geometry{margin=2.2cm}%
%
\usepackage{amsmath}%
\usepackage{amssymb}%
\usepackage{booktabs}%
\usepackage{longtable}%
\usepackage{array}%
\usepackage{hyperref}%
\hypersetup{hidelinks}%
\title{Soluções\\[4pt]\large Comparação com o gabarito Gabarito preliminar oficial (SSPM, 01/09/2026), prova AMARELA\thanks{Gerado com gpt-resolve: \url{https://github.com/lgabs/gpt-resolve}}}%
\date{29/10/2026}%
%
\begin{document}%
\normalsize%
\maketitle%
\section*{Resumo}\textbf{Nota média: 10,0/10} (22 de 22 itens com a resolução correta; cada item vale o mesmo dentro da questão).\par\medskip%
\begin{longtable}{@{}c c p{0.78\textwidth}@{}}\toprule Questão & Nota & Resumo \\ \midrule
1 & 10,0 & A resolução e o gabarito preliminar oficial (prova Amarela) coincidem na alternativa D, e as duas estão corretas. I, III e IV são verdadeiras, e II é falsa porque f(x) = x$\cdot$e\^{}\{1/x\} tem duas assíntotas, x = 0 e y = x + 1. Refiz cada afirmativa e conferi numericamente os limites, a monotonia de a\_n e a desigualdade do item IV. A argumentação da resolução é rigorosa e não encontrei erro. Ficaram só duas ressalvas leves: o caso $\gamma$ = $\beta$ na afirmativa I e a forma de contar a oblíqua na II, e nenhuma delas muda a resposta. \\
2 & 10,0 & A resolução está matematicamente correta. As raízes reais são (3-$\surd$5)/2, 1 e (3+$\surd$5)/2, e a razão menor/maior é (7-3$\surd$5)/2 $\approx$ 0,146, valor que não aparece em nenhuma alternativa. O gabarito preliminar oficial marca a (A), (3-$\surd$5)/2 $\approx$ 0,382. Esse número é a razão da PG das raízes, não a razão entre as raízes extremas, então o gabarito está errado pelo enunciado literal e a questão deveria ser anulada. A resolução aponta isso e indica a (A) só como escolha forçada, que coincide com o gabarito. \\
3 & 10,0 & A resolução e o gabarito preliminar oficial coincidem na alternativa (C), arccos(-$\surd$5/3), e os dois estão certos. Conferi o ângulo diedro do icosaedro regular de forma independente, pela lei dos cossenos e por coordenadas: cos $\theta$ = -$\surd$5/3, ou seja, $\theta$ $\approx$ 138,19$^\circ$. A resolução descarta as demais alternativas corretamente: (A) e (B) dão ângulos negativos, e (D) e (E) estão fora do domínio da função. Não há erro de conta nem de raciocínio; ficam só ressalvas menores de interpretação e de rigor. \\
4 & 10,0 & A resolução e o gabarito preliminar oficial (prova Amarela) concordam: a resposta é a alternativa B (Laplace e Cauchy), e está correta. A condição (1), com i = k, é a expansão de Laplace pela linha i. A condição (2), com i $\neq$ k, é o Teorema de Cauchy: a soma dos elementos de uma linha multiplicados pelos cofatores de outra linha paralela é zero. A justificativa da resolução é completa e rigorosa, e a conferência para n = 2 bate. Não há nenhuma divergência para levar ao professor. \\
5 & 10,0 & A resolução está correta e coincide com o gabarito preliminar oficial: alternativa (E), 27$\pi$. As interseções são x = 2 e x = 4, e nesse intervalo x$^2$ > 2\^{}x. O volume é V = $\pi$(992/5 - 120/ln 2), que dá $\approx$ 26,97$\pi$ com o dado ln 2 $\cong$ 0,7. Conferi todas as contas numericamente. O único ponto a observar: com o ln 2 exato o volume seria $\approx$ 25,28$\pi$, mais perto de 24$\pi$. O enunciado, porém, fixa ln 2 $\cong$ 0,7, então (E) é a resposta pretendida. \\
6 & 10,0 & Questão 6: a resolução marca (B) 5$\pi$/3 e o gabarito preliminar oficial também dá B, então as respostas coincidem e as duas estão certas. Pela identidade arcsen u + arccos u = $\pi$/2, f é constante e vale $\pi$/2 no domínio {[}{-}1/3, 3{]}. A região é um retângulo de base 10/3 e altura $\pi$/2, com área 5$\pi$/3. Conferi o domínio e a conta nos extremos, e a transcrição do enunciado está fiel ao PDF original. Não encontrei nenhum ponto relevante para levar ao professor. \\
7 & 10,0 & A resolução e o gabarito oficial preliminar concordam: a resposta é (A), m < -4 ou m > 4. Refiz a conta e cheguei ao mesmo resultado: a curva é hipérbole se e somente se det{[}{[}1,m{]},{[}m,16{]}{]} = 16 - m$^2$ < 0, e ela não degenera porque o termo independente não é nulo. A resolução está correta e completa, inclusive na análise de m = $\pm$4 (par de retas paralelas) e de |m| < 4 (elipse). Não encontrei erros. \\
8 & 10,0 & A resolução está correta e bate com o gabarito oficial preliminar: alternativa (A), 12/$\surd$154. Os vetores diretores foram lidos corretamente, as contas estão certas e a observação de que as retas são reversas é verdadeira. Usar o módulo no numerador é o certo, e aqui não muda nada, porque o produto escalar já é positivo. Não há pontos a levar ao professor. \\
9 & 10,0 & A resolução e o gabarito concordam: a resposta é a alternativa (A). Para x $\neq$ $\pm$1, a função se simplifica para 1/(1-x). O gráfico tem assíntotas x = 1 e y = 0 e uma descontinuidade removível (buraco) em (-1, 1/2). Por causa desse ponto excluído, a resposta é (A) e não (E). Conferi o enunciado no PDF original e verifiquei os valores e limites numericamente. Não encontrei erro. \\
10 & 10,0 & A resolução está correta e coincide com o gabarito: alternativa (E), V, F, V, V. Na ordem usual (V,V), (V,F), (F,V), (F,F), a tabela{-}verdade de \textasciitilde{}n $\to$ \textasciitilde{}m é igual à de m $\to$ n e só é falsa quando m = V e n = F. A justificativa linha a linha está certa e a resolução declara a ordem que usa. O único ponto a observar é que o enunciado não diz essa ordem, mas a ordem padrão é a única que leva a uma das alternativas sem ambiguidade. \\
11 & 10,0 & A resolução e o gabarito preliminar oficial (prova AMARELA) concordam: a alternativa é a (B), e\^{}\{{-}1\}, e está correta. Com u = tg($\pi$x/4), o limite vira u\^{}\{2u/(1-u$^2$)\}. O expoente tende a -1 pelo limite fundamental de ln u/(u-1), então o resultado é e\^{}\{{-}1\}. A checagem numérica em Python dá $\approx$ 0,36788 pelos dois lados. A transcrição do enunciado confere com o PDF original. \\
12 & 10,0 & A resolução está correta. Em {[}0;5) as curvas se cruzam em x=2 e x=4, com x$^2$ $\geq$ 2\^{}x entre elas, e pelo método das arruelas V = $\pi$(198,4 - 120/ln 2). Com o dado ln 2 $\cong$ 0,7, V $\cong$ 26,97$\pi$ $\approx$ 27$\pi$, alternativa (E). O gabarito preliminar oficial anulou a questão e aceita qualquer alternativa, então a resposta (E) recebe o ponto. A anulação também se justifica: o valor exato, $\approx$25,28$\pi$, não corresponde a nenhuma alternativa, o resultado depende muito da aproximação de ln 2 e o 'x$\in${[}0;5)' deixa a região ambígua. A própria resolução aponta essa sensibilidade. \\
13 & 10,0 & A resolução está correta e bate com o gabarito preliminar oficial: alternativa (C), $\alpha$$^2$. Derivando xy = 4 duas vezes com ? = $\alpha$ e ÿ = 0, no ponto (2,2) obtém{-}se ? = -$\alpha$ e ? = $\alpha$$^2$. Conferi a conta à mão e numericamente. Ficam apenas duas observações menores: a leitura do enunciado e a redação da análise dimensional. \\
14 & 10,0 & A resolução e o gabarito oficial preliminar (SSPM, prova Amarela) concordam: a resposta é a alternativa (B), 128/15. Refiz a conta de forma independente, com frações exatas e com integração numérica, e o resultado foi o mesmo. A interseção, o sinal de g {-} f e a primitiva estão todos certos. A resolução está correta e completa, sem erros a apontar. \\
15 & 10,0 & A resolução e o gabarito preliminar oficial (B) concordam, e os dois estão certos: a resposta é 3/5. No primeiro turno, P(A)=3/11 e P(D)=2/11. No segundo turno vale a razão 3:2, o que dá P(A)=3/5. A transcrição confere com o original. Há só uma ressalva menor: o enunciado supõe que um candidato vence no primeiro turno, mas isso não afeta a resposta, que depende só da razão entre as chances de A e D. \\
16 & 10,0 & Resolução e gabarito oficial preliminar dão a mesma resposta: (D) 81/16. A soma da série é 1/2 + 1 = 3/2, e (3/2)\^{}4 = 81/16. A transcrição confere com o original, todas as contas da resolução, inclusive as da verificação, estão corretas e não há nada a levar ao professor. \\
17 & 10,0 & A resolução e o gabarito oficial preliminar concordam: a alternativa é a (A). A resolução está certa. A relação de Stifel vale para x real e faz a soma se cancelar em cadeia até sobrar ({-}1)\^{}n$\cdot$C(x{-}1,n) = ({-}1)\^{}n(x{-}1)$\cdot$$\cdot$$\cdot$(x{-}n)/n!. Conferi a identidade com frações exatas e ela vale. As outras alternativas foram eliminadas com argumentos corretos. Não encontrei erro nem na resolução nem no gabarito. \\
18 & 10,0 & A resolução e o gabarito oficial concordam: a resposta é (A) a = 1/31. A dedução está correta: interseção em x$^2$ = 8/(a+1) e ortogonalidade com -4ax$^2$ = -1, o que dá a = 1/31. A verificação também confere: os pontos são x = $\pm$$\surd$31/2 e y = -27/4. As contas das outras alternativas também estão certas. Não há erro na resolução nem no gabarito. \\
19 & 10,0 & A resolução e o gabarito oficial preliminar marcam a mesma alternativa, (A), e os dois estão certos. A primitiva F(x) = -x$\cdot$cot x + ln(sen x) - x$^2$/2 foi conferida por derivação. Com ela, I = -ln($\surd$2/2) + $\pi$/4 - 3$\pi$$^2$/32 $\approx$ 0,2067, e esse valor bate com a integração numérica. As demais alternativas têm valores bem diferentes do da integral. O enunciado transcrito confere com o original. \\
20 & 10,0 & A resolução e o gabarito preliminar oficial (prova Amarela) marcam a mesma alternativa, (D) R\$ 132.300,00, e os dois estão certos: 150.000 $\times$ 0,9 $\times$ 0,7 $\times$ 1,4 = 132.300. A resolução declara como entende o enunciado e usa corretamente os fatores multiplicativos para variações sucessivas. O único ponto a observar é a redação: o enunciado chama 150.000 de valor que o veículo \textbackslash{}"custa atualmente\textbackslash{}", mas a conta só fecha com uma alternativa se 150.000 for o valor de aquisição. \\
21 & 10,0 & A resolução e o gabarito oficial preliminar marcam a mesma alternativa, (E), e os dois estão certos. Como f(x)=cos(2$\cdot$arccos x)=2x$^2$-1, temos g=1/(2x$^2$-1). Esse gráfico tem assíntotas verticais em x=$\pm$$\surd$2/2, um ramo central negativo com máximo -1 em x=0 e ramos laterais positivos que descem de +$\infty$ até 1 em x=$\pm$1, e é exatamente o que mostra o gráfico (E). A resolução está completa e correta. \\
22 & 10,0 & A resolução e o gabarito concordam: alternativa (D), (1+$\surd$2)/8. Refiz tudo de forma independente: I = BE?CH = (-4$\surd$2, 0), área de IEG = 16(1+$\surd$2)$^2$ e área do octógono = 128(1+$\surd$2), o que dá razão (1+$\surd$2)/8 $\approx$ 0,302. A resolução está correta e bem justificada, e ainda confere a área de IEG pelo produto vetorial. Ficam só observações menores: a figura não está em escala e a interpretação da colinearidade de C, I e H não foi declarada. \\
\bottomrule\end{longtable}%
\section*{Principais pontos de dúvida}%
\begin{itemize}%
\item%
\textbf{Q2} \textit{(key)} O gabarito preliminar (A) não corresponde ao que o enunciado pede. A razão menor/maior raiz real é (7-3$\surd$5)/2 $\approx$ 0,146, e não (3-$\surd$5)/2 $\approx$ 0,382. Nenhuma alternativa traz esse valor, então a questão é forte candidata a anulação. Vale conferir no gabarito definitivo se a banca anulou ou manteve a (A).%
\end{itemize}%
\section*{Questão 1}%
\begin{longtable}{@{}p{0.06\textwidth} p{0.33\textwidth} p{0.33\textwidth} p{0.1\textwidth} p{0.08\textwidth}@{}}\toprule Item & Resolução & Gabarito & Veredito & Certo \\ \midrule
1 & (D) Apenas as afirmativas I, III e IV são verdadeiras & D & igual & ambos \\
\bottomrule\end{longtable}%
\begin{itemize}%
\item%
\textbf{1)} Conferi o enunciado no PDF (página 2/12, prova Amarela) e a transcrição está fiel. Refiz as quatro afirmativas. I é verdadeira: $\alpha$ não pode ser paralelo a $\gamma$ ? $\beta$. Como r ? $\beta$ e s ? $\gamma$, com $\beta$ ? $\gamma$ = ?, as retas r e s, que estão ambas em $\alpha$, não têm ponto comum, logo são paralelas. II é falsa: f(x) = x$\cdot$e\^{}\{1/x\} tem assíntota vertical x = 0, porque f $\to$ +$\infty$ quando x $\to$ 0?, e assíntota oblíqua y = x + 1 nos dois sentidos. Pelo cálculo numérico, f(x) - x $\to$ 1 para x = $\pm$10$^{5}$ e f(0,01) $\approx$ 2,7$\cdot$10$^{4}$$^{1}$. São duas assíntotas, não apenas uma. III é verdadeira: pela desigualdade MA--MG estrita aplicada a n cópias de (1 + 1/n) e ao número 1, vale a\_n < a\_\{n+1\}. Numericamente, a sequência dá 2; 2,25; 2,370; 2,441; \ldots{} IV é verdadeira: pela MA--MG em $\surd$x, $\surd$y e $\surd$z, elevando ao quadrado. Num teste com 10$^{5}$ triplas aleatórias, a diferença mínima foi $\geq$ 0. A resolução, portanto, chega à alternativa D, que coincide com o gabarito, e a argumentação está correta e completa.%
\end{itemize}%
\paragraph{Pontos de dúvida}%
\begin{itemize}%
\item%
\textit{(statement)} Na afirmativa I, 'qualquer plano paralelo a $\beta$' pode incluir o próprio $\beta$, se planos coincidentes contarem como paralelos. Nesse caso, a interseção é a própria r, que também é paralela a r pela convenção de paralelismo amplo. A resolução supõe $\gamma$ $\neq$ $\beta$ sem tratar esse caso, mas o valor lógico de I não muda.%
\item%
\textit{(statement)} Na afirmativa II, contar a oblíqua y = x + 1 uma vez só ou uma vez para cada sentido (x $\to$ +$\infty$ e x $\to$ -$\infty$) não altera nada: em qualquer contagem há pelo menos duas assíntotas (a vertical x = 0 mais a oblíqua), e II continua falsa.%
\end{itemize}%
\section*{Questão 2}%
\begin{longtable}{@{}p{0.06\textwidth} p{0.33\textwidth} p{0.33\textwidth} p{0.1\textwidth} p{0.08\textwidth}@{}}\toprule Item & Resolução & Gabarito & Veredito & Certo \\ \midrule
Q2 & Razão menor/maior raiz real = (7-3$\surd$5)/2 $\approx$ 0,146. Nenhuma alternativa confere, então a questão seria passível de anulação. Se for obrigatório marcar uma, a resolução marca a (A) (3-$\surd$5)/2. & A, (3-$\surd$5)/2 $\approx$ 0,382 & diferente & resolução \\
\bottomrule\end{longtable}%
\begin{itemize}%
\item%
\textbf{Q2)} Refiz a conta. (x-1)(x$^{4}$-4x$^3$+5x$^2$-4x+1) = x$^{5}$-5x$^{4}$+9x$^3$-9x$^2$+5x-1, e x$^{4}$-4x$^3$+5x$^2$-4x+1 = (x$^2$-x+1)(x$^2$-3x+1). As expansões conferem. O fator x$^2$-x+1 tem $\Delta$=-3, então só há três raízes reais: (3-$\surd$5)/2 $\approx$ 0,382; 1; e (3+$\surd$5)/2 $\approx$ 2,618. Conferi numericamente que P se anula nesses três pontos. A razão menor/maior é (3-$\surd$5)/(3+$\surd$5) = (7-3$\surd$5)/2 $\approx$ 0,1459. Esse valor não está entre as opções. A (A), que é o gabarito, vale $\approx$ 0,382: é a própria menor raiz, ou a razão da PG (3+$\surd$5)/2, 1, (3-$\surd$5)/2. A razão pedida é o quadrado desse número. A (A) só sairia se a banca tivesse dividido a menor raiz por 1, ou confundido a razão entre as raízes extremas com a razão da PG. Pelo enunciado literal, o gabarito está errado e a questão não tem alternativa correta. A resolução identifica isso, dá o valor certo e indica a (A) apenas como escolha forçada. Por isso a opção marcada coincide com o gabarito, mas a conclusão matemática é diferente, e quem está certo é a resolução.%
\end{itemize}%
\paragraph{Pontos de dúvida}%
\begin{itemize}%
\item%
\textit{(key)} O gabarito preliminar (A) não corresponde ao que o enunciado pede. A razão menor/maior raiz real é (7-3$\surd$5)/2 $\approx$ 0,146, e não (3-$\surd$5)/2 $\approx$ 0,382. Nenhuma alternativa traz esse valor, então a questão é forte candidata a anulação. Vale conferir no gabarito definitivo se a banca anulou ou manteve a (A).%
\item%
\textit{(statement)} Conferi a transcrição com o PDF original (prova Amarela, pág. 2/12) e ela está fiel. O enunciado diz claramente 'razão entre a menor raiz real e a maior raiz real', ou seja, o quociente. Pela forma da alternativa (A), a banca provavelmente pensou na razão da PG formada pelas raízes reais.%
\item%
\textit{(solution)} A resolução dá duas respostas: o valor exato fora das alternativas e a (A) como opção forçada. Em uma correção que exija só a letra marcada, ela 'acertaria' a (A), mas pelo motivo errado em relação ao enunciado. O critério de pontuação deve ser decidido pelo professor.%
\end{itemize}%
\section*{Questão 3}%
\begin{longtable}{@{}p{0.06\textwidth} p{0.33\textwidth} p{0.33\textwidth} p{0.1\textwidth} p{0.08\textwidth}@{}}\toprule Item & Resolução & Gabarito & Veredito & Certo \\ \midrule
Questão 3 (única) & (C) arccos(-$\surd$5/3) $\approx$ 138,19$^\circ$ & C & igual & ambos \\
\bottomrule\end{longtable}%
\begin{itemize}%
\item%
\textbf{Questão 3 (única))} O enunciado da resolução confere com o PDF original (prova Amarela, pág. 2/12). Refiz a conta de dois jeitos. (1) Lei dos cossenos no triângulo BME: BE$^2$ = $\varphi$$^2$a$^2$ = (3+$\surd$5)a$^2$/2 e BM = EM = a$\surd$3/2, o que dá cos $\theta$ = -$\surd$5/3. (2) Coordenadas dos vértices (0, $\pm$1, $\pm$$\varphi$) e permutações cíclicas: o produto escalar normalizado das alturas que partem do ponto médio da aresta comum vale -0,745356 = -$\surd$5/3, e $\theta$ $\approx$ 138,19$^\circ$. A análise das alternativas também está certa. Em (A) e (B), os valores ficam em $\approx$ -18,63$^\circ$ e $\approx$ -29,21$^\circ$, negativos e fora de (0, $\pi$). Em (D), cosec $\theta$ = -$\surd$5/5 exigiria |sen $\theta$| = $\surd$5 > 1. Em (E), |-$\surd$5/2| > 1 sai do domínio do arccos. Os valores auxiliares sen $\theta$ = 2/3 e tg $\theta$ = -2$\surd$5/5 estão corretos.%
\end{itemize}%
\paragraph{Pontos de dúvida}%
\begin{itemize}%
\item%
\textit{(statement)} 'Ângulo entre duas faces adjacentes' foi lido como o ângulo diedro interno, que é obtuso ($\approx$138,19$^\circ$). Se fosse lido como o ângulo entre os planos (agudo), o resultado seria arccos($\surd$5/3) $\approx$ 41,81$^\circ$, que não está entre as alternativas. Só a leitura como diedro tem alternativa correspondente, então isso não afeta a nota.%
\item%
\textit{(solution)} No passo 1, a resolução diz que os 5 vizinhos de um vértice são coplanares e formam um pentágono regular, e justifica pela simetria de rotação de 72$^\circ$. O argumento é válido, mas usa sem demonstrar que essa rotação é uma simetria do icosaedro regular. É um fato padrão e aceitável.%
\end{itemize}%
\section*{Questão 4}%
\begin{longtable}{@{}p{0.06\textwidth} p{0.33\textwidth} p{0.33\textwidth} p{0.1\textwidth} p{0.08\textwidth}@{}}\toprule Item & Resolução & Gabarito & Veredito & Certo \\ \midrule
única & (B) Laplace e Cauchy & B & igual & ambos \\
\bottomrule\end{longtable}%
\begin{itemize}%
\item%
\textbf{única)} Conferi o enunciado no PDF (página 2/12, prova Amarela) e a transcrição está fiel. Para i = k, b\_ii = $\Sigma$\_j a\_ij$\cdot$A\_ij é a expansão do determinante pela linha i, ou seja, o Teorema de Laplace. Para i $\neq$ k, b\_ik = $\Sigma$\_j a\_ij$\cdot$A\_kj multiplica os elementos de uma linha pelos cofatores de outra linha, paralela a ela, e dá zero. Nos livros brasileiros (Iezzi, Fundamentos vol. 4) esse é o Teorema de Cauchy. A ordem (1) Laplace e (2) Cauchy é a alternativa B. A demonstração da resolução está correta: troca a linha k por uma cópia da linha i, os cofatores da linha k não mudam e a nova matriz tem determinante nulo. A conferência para n = 2 também está certa: ad - bc e -ab + ba = 0. A alternativa D traz os mesmos teoremas na ordem trocada. Chió é um método para abaixar a ordem do determinante, Binet é det(AB) = det A$\cdot$det B e Cramer resolve sistemas lineares. Nenhum desses serve aqui.%
\end{itemize}%
\paragraph{Pontos de dúvida}%
\begin{itemize}%
\item%
\textit{(statement)} O enunciado usa B para a matriz, a\_ij para os elementos de B e B\_jk para os elementos da adjunta. Também usa b\_ik para os elementos do produto e escreve 'Kramer' em vez de Cramer. Essa notação misturada não muda a resposta, e a resolução declarou a convenção (B\_jk = A\_kj) de forma explícita.%
\item%
\textit{(key)} O nome 'Teorema de Cauchy' para a soma dos elementos de uma fila multiplicados pelos cofatores de uma fila paralela é uma convenção dos livros didáticos brasileiros. Fora do Brasil esse resultado costuma ser tratado como parte da expansão de Laplace ('alien cofactors'). Dentro das alternativas dadas, a B é a única coerente.%
\end{itemize}%
\section*{Questão 5}%
\begin{longtable}{@{}p{0.06\textwidth} p{0.33\textwidth} p{0.33\textwidth} p{0.1\textwidth} p{0.08\textwidth}@{}}\toprule Item & Resolução & Gabarito & Veredito & Certo \\ \midrule
único & (E) 27$\pi$. V = $\pi$$\int$$_{2}$$^{4}$(x$^{4}$ - 4\^{}x)dx = $\pi$(992/5 - 120/ln 2) $\cong$ 26,97$\pi$ com ln 2 $\cong$ 0,7 & E & igual & ambos \\
\bottomrule\end{longtable}%
\begin{itemize}%
\item%
\textbf{único)} Conferi o enunciado no PDF original (pág. 3 do arquivo, Questão 5) e ele bate com a transcrição. Em {[}0;5) as curvas se cruzam só em x = 2 e x = 4; a raiz negativa, x $\approx$ -0,767, fica fora do intervalo. Em (2;4) vale x$^2$ > 2\^{}x, então a única região limitada pelas interseções é 2 $\leq$ x $\leq$ 4, 2\^{}x $\leq$ y $\leq$ x$^2$. Pelo método das arruelas, V = $\pi$$\int$$_{2}$$^{4}$(x$^{4}$ - 4\^{}x)dx = $\pi${[}(4$^{5}$ - 2$^{5}$)/5 - (4$^{4}$ - 4$^2$)/ln 4{]} = $\pi$(198,4 - 120/ln 2). Com ln 2 = 0,7 dá V/$\pi$ = 26,971 $\approx$ 27, alternativa (E). Refiz a conta em Python: a fórmula fechada e a regra de Simpson dão V/$\pi$ = 25,2766 com o ln 2 exato, o que confirma a integral. Resolução e gabarito estão certos e concordam.%
\end{itemize}%
\paragraph{Pontos de dúvida}%
\begin{itemize}%
\item%
\textit{(statement)} A resposta (E) depende de usar ln 2 $\cong$ 0,7, como o enunciado manda. Com o valor exato, ln 2 = 0,6931\ldots{}, o volume é $\approx$ 25,28$\pi$. Esse valor fica mais perto de 24$\pi$ (D), a 1,28 de distância, do que de 27$\pi$ (E), a 1,72. O resultado é a diferença de dois números grandes, 198,4 e $\approx$ 173, e por isso é muito sensível à aproximação de ln 2. A resolução aponta isso de forma correta. Como o dado é explícito, (E) é a resposta pretendida, mas vale citar o caso numa eventual discussão de recurso.%
\item%
\textit{(statement)} O enunciado diz 'x $\in$ {[}0;5)', mas a única região fechada pelas interseções é {[}2;4{]}. As faixas {[}0;2{]} e {[}4;5) não são limitadas pelas interseções nas duas pontas. Se alguém interpretasse que o sólido vem da região entre as curvas em todo o intervalo {[}0;5), o volume seria muito maior ($\approx$ 163$\pi$ com ln 2 exato) e não bateria com nenhuma alternativa. Isso reforça a leitura adotada pela resolução e pelo gabarito.%
\end{itemize}%
\section*{Questão 6}%
\begin{longtable}{@{}p{0.06\textwidth} p{0.33\textwidth} p{0.33\textwidth} p{0.1\textwidth} p{0.08\textwidth}@{}}\toprule Item & Resolução & Gabarito & Veredito & Certo \\ \midrule
6 & (B) 5$\pi$/3 & B & igual & ambos \\
\bottomrule\end{longtable}%
\begin{itemize}%
\item%
\textbf{6)} Conferi a transcrição com o PDF original (página 3/12, prova Amarela). O enunciado e as alternativas batem. Pela identidade arcsen u + arccos u = $\pi$/2, válida para u $\in$ {[}{-}1,1{]}, f vale $\pi$/2 em todo o domínio. A condição {-}1 $\leq$ (2x{-}1)/(x+2) $\leq$ 1 dá (x{-}3)/(x+2) $\leq$ 0 e (3x+1)/(x+2) $\geq$ 0. Daí vêm {-}2 < x $\leq$ 3 e (x < {-}2 ou x $\geq$ {-}1/3). A interseção é D = {[}{-}1/3, 3{]}. Nos extremos, u vale {-}1 em x = {-}1/3 e 1 em x = 3. A área é a de um retângulo: (3 + 1/3)$\cdot$($\pi$/2) = (10/3)($\pi$/2) = 5$\pi$/3 $\approx$ 5,236, que é a alternativa B. A resolução e o gabarito coincidem e estão corretos.%
\end{itemize}%
\paragraph{Pontos de dúvida}%
\begin{itemize}%
\item%
\textit{(solution)} Na conclusão, a resolução diz que a alternativa (A) 'se obtém ao esquecer o trecho {[}{-}1/3, 0{]}'. Isso é só uma especulação sobre como alguém chegaria à alternativa errada: a base 3 corresponde à base usada em (A), mas não há como saber se esse é o erro típico. Também não aparece de onde viriam as bases das alternativas (C) e (D). Nada disso afeta a resposta.%
\end{itemize}%
\section*{Questão 7}%
\begin{longtable}{@{}p{0.06\textwidth} p{0.33\textwidth} p{0.33\textwidth} p{0.1\textwidth} p{0.08\textwidth}@{}}\toprule Item & Resolução & Gabarito & Veredito & Certo \\ \midrule
Q7 & (A) m < {-}4 ou m > 4 & A & igual & ambos \\
\bottomrule\end{longtable}%
\begin{itemize}%
\item%
\textbf{Q7)} Conferi o enunciado no PDF original (página 3/12 da prova Amarela): a transcrição está fiel. Refiz a conta. A forma quadrática tem matriz M = {[}{[}1,m{]},{[}m,16{]}{]}, com det M = 16 - m$^2$ (o mesmo que B$^2$ - 4AC = 4m$^2$ - 64 > 0). Os termos lineares são nulos e o segundo membro é 1 $\neq$ 0, então a equação vira $\lambda$$_{1}$X$^2$ + $\lambda$$_{2}$Y$^2$ = 1. Ela é hipérbole exatamente quando $\lambda$$_{1}$$\lambda$$_{2}$ < 0, ou seja, |m| > 4. O determinante 3$\times$3 da cônica, m$^2$ - 16, é não nulo nesse intervalo, o que confirma que ela não é degenerada. Para |m| < 4 a curva é uma elipse real, porque o traço vale 17 > 0. Para m = $\pm$4 temos (x $\pm$ 4y)$^2$ = 1, um par de retas paralelas. As alternativas (D) e (E) são partes do conjunto, e (B) e (C) incluem valores que dão elipse. A resolução e o gabarito dão a mesma resposta, (A), e os dois estão certos.%
\end{itemize}%
\paragraph{Pontos de dúvida}%
\begin{itemize}%
\item%
\textit{(statement)} A expressão "valor estabelecido plenamente" é pouco clara. A leitura natural é "o conjunto completo dos valores de m", e é ela que separa (A) de (D) e (E), que também só dão hipérboles mas não cobrem o conjunto todo. A resolução adota essa leitura de forma explícita.%
\end{itemize}%
\section*{Questão 8}%
\begin{longtable}{@{}p{0.06\textwidth} p{0.33\textwidth} p{0.33\textwidth} p{0.1\textwidth} p{0.08\textwidth}@{}}\toprule Item & Resolução & Gabarito & Veredito & Certo \\ \midrule
único & (A) 12/$\surd$154 & A & igual & ambos \\
\bottomrule\end{longtable}%
\begin{itemize}%
\item%
\textbf{único)} Conferi o enunciado original (pág. 3/12 do PDF, prova Amarela) e ele bate com a transcrição. Os vetores diretores são u=(1,2,3) e v=(1,1,3). O produto escalar é u$\cdot$v=12, com |u|=$\surd$14 e |v|=$\surd$11, então cos $\theta$=12/$\surd$154$\approx$0,967 ($\theta$$\approx$14,8$^\circ$), conferido em python. A verificação extra de que as retas são reversas também está certa: igualando as duas primeiras coordenadas sai a=t=0, e a terceira coordenada dá 0=3, o que é absurdo. A forma 6$\surd$154/77 é equivalente.%
\end{itemize}%
\section*{Questão 9}%
\begin{longtable}{@{}p{0.06\textwidth} p{0.33\textwidth} p{0.33\textwidth} p{0.1\textwidth} p{0.08\textwidth}@{}}\toprule Item & Resolução & Gabarito & Veredito & Certo \\ \midrule
Q9 & (A): para x $\neq$ $\pm$1, f(x) = 1/(1-x). Assíntotas x = 1 e y = 0, ramo esquerdo acima do eixo x e crescente, ramo direito abaixo do eixo x e crescente, buraco em (-1, 1/2). & A & igual & ambos \\
\bottomrule\end{longtable}%
\begin{itemize}%
\item%
\textbf{Q9)} Conferi o enunciado no PDF original (pág. 4/12 da prova Amarela), e ele coincide com a transcrição. Para x $\neq$ -1, f(x) = -(x+1)/((x-1)(x+1)) = 1/(1-x). Verifiquei numericamente: f $\to$ 1/2 quando x $\to$ -1 (buraco em (-1, 1/2), à esquerda do eixo y e acima do eixo x) e f(0) = 1. Quando x $\to$ 1$^{-}$, f $\to$ +$\infty$, e quando x $\to$ 1?, f $\to$ -$\infty$. Para x $\to$ $\pm$$\infty$, f $\to$ 0? pela esquerda e 0$^{-}$ pela direita. O gráfico (A) tem exatamente isso: assíntota à direita do eixo y, ramo esquerdo positivo e crescente com a bolinha aberta em x < 0, ramo direito negativo e crescente. O gráfico (E) tem o mesmo formato, mas não mostra o ponto excluído. O (D) corresponde a -1/(x+1), com buraco em (1, -1/2). Os gráficos (B) e (C) têm sinal e monotonicidade invertidos. A resolução está correta e completa.%
\end{itemize}%
\paragraph{Pontos de dúvida}%
\begin{itemize}%
\item%
\textit{(statement)} Os esboços não têm escala, então não dá para confirmar que a bolinha de (A) está exatamente em (-1, 1/2). Dá para checar só a posição qualitativa: x < 0 e y > 0. Isso basta para decidir entre (A) e (E).%
\end{itemize}%
\section*{Questão 10}%
\begin{longtable}{@{}p{0.06\textwidth} p{0.33\textwidth} p{0.33\textwidth} p{0.1\textwidth} p{0.08\textwidth}@{}}\toprule Item & Resolução & Gabarito & Veredito & Certo \\ \midrule
10 & (E) V, F, V, V & E & igual & ambos \\
\bottomrule\end{longtable}%
\begin{itemize}%
\item%
\textbf{10)} A transcrição confere com o original (página 4/12 do PDF, prova Amarela). Na ordem usual das linhas, (m,n) = (V,V), (V,F), (F,V), (F,F), a coluna de \textasciitilde{}n $\to$ \textasciitilde{}m fica V, F, V, V. Ela é igual à coluna de m $\to$ n, porque a contrapositiva é equivalente à condicional. A sentença só é falsa quando m = V e n = F. Conferi as quatro linhas da tabela{-}verdade e a justificativa de cada uma, e estão corretas. Só as alternativas (A) e (E) têm exatamente um F, como deve ter uma condicional. A (E) põe esse F na posição da linha (V,F) na ordem padrão. A resolução e o gabarito coincidem e estão certos.%
\end{itemize}%
\paragraph{Pontos de dúvida}%
\begin{itemize}%
\item%
\textit{(statement)} O enunciado não diz em que ordem estão as linhas da tabela{-}verdade, e a resposta depende disso. A resolução declara que usa a ordem usual (V,V), (V,F), (F,V), (F,F). Outras ordens comuns, como começar por (F,F) ou trocar os papéis de m e n, dão V, V, F, V, que não está entre as alternativas. A alternativa (A), com o F no fim, exigiria uma ordem incomum que deixasse a linha (m=V, n=F) por último. Por isso a (E) é a única resposta coerente com a convenção padrão, e essa omissão não deve mudar a resposta.%
\end{itemize}%
\section*{Questão 11}%
\begin{longtable}{@{}p{0.06\textwidth} p{0.33\textwidth} p{0.33\textwidth} p{0.1\textwidth} p{0.08\textwidth}@{}}\toprule Item & Resolução & Gabarito & Veredito & Certo \\ \midrule
11 & (B) e\^{}\{{-}1\} & B & igual & ambos \\
\bottomrule\end{longtable}%
\begin{itemize}%
\item%
\textbf{11)} O enunciado confere com o original (pág. 3 do PDF, 'Página 4/12'): lim\_\{x$\to$1\} (tg $\pi$x/4)\^{}(tg $\pi$x/2). Refiz a conta com u = tg($\pi$x/4). Pelo arco duplo, tg($\pi$x/2) = 2u/(1-u$^2$). O expoente do logaritmo fica 2u$\cdot$ln u/(1-u$^2$) = -{[}2u/(u+1){]}$\cdot${[}ln u/(u-1){]}, que tende a -1$\cdot$1 = -1. Logo L = e\^{}\{{-}1\}. Conferi numericamente em Python, com x = 1 $\pm$ 1e{-}4 e x = 1 $\pm$ 1e{-}7: deu $\approx$ 0,3678794 pelos dois lados, igual a 1/e $\approx$ 0,3678794. A resolução e o gabarito marcam a mesma alternativa (B), e ela está certa.%
\end{itemize}%
\paragraph{Pontos de dúvida}%
\begin{itemize}%
\item%
\textit{(solution)} O parágrafo 'Verificação dos limites laterais' é só uma checagem qualitativa: diz que a potência 'tende a um valor menor que 1' e não prova nada. Isso não afeta o resultado, porque a conta principal com u$\to$1 já vale pelos dois lados. É só uma questão de redação.%
\end{itemize}%
\section*{Questão 12}%
\begin{longtable}{@{}p{0.06\textwidth} p{0.33\textwidth} p{0.33\textwidth} p{0.1\textwidth} p{0.08\textwidth}@{}}\toprule Item & Resolução & Gabarito & Veredito & Certo \\ \midrule
12 & (E) 27$\pi$. V = $\pi$$\int$$_{2}$$^{4}$(x$^{4}$ - 4\^{}x)dx = $\pi$(198,4 - 120/ln 2) $\cong$ 26,97$\pi$ com ln 2 $\cong$ 0,7 & Questão anulada pela banca (qualquer alternativa é aceita) & igual & ambos \\
\bottomrule\end{longtable}%
\begin{itemize}%
\item%
\textbf{12)} Conferi o enunciado no PDF (pág. 4/12, prova Amarela) e ele coincide com a transcrição: f(x)=2\^{}x, g(x)=x$^2$, x$\in${[}0;5), Dado ln 2 $\cong$ 0,7, alternativas de 17$\pi$ a 27$\pi$. Refiz a conta. Em {[}0;5) as curvas só se cruzam em x=2 e x=4 (a raiz negativa, $\approx$ -0,767, fica fora), e em {[}2,4{]} vale x$^2$ $\geq$ 2\^{}x. Pelo método das arruelas, $\int$$_{2}$$^{4}$x$^{4}$dx = 198,4 e $\int$$_{2}$$^{4}$4\^{}x dx = 120/ln 2. Com ln 2 = 0,7, V = (198,4 - 171,43)$\pi$ = 26,971$\pi$ $\approx$ 27$\pi$, que é a alternativa (E). Conferi por integração numérica (Simpson): com o ln 2 exato, V = 25,2766$\pi$. A resolução está correta e usa o dado do enunciado. Como a banca anulou a questão e aceita qualquer alternativa, a resposta (E) recebe o ponto. Por isso marquei 'match': a resposta é compatível com o gabarito e não perde nota. A anulação também se justifica. O valor exato, $\approx$25,3$\pi$, não bate com nenhuma alternativa. O resultado depende muito da aproximação de ln 2, porque se subtraem números próximos (198,4 - 171,4 contra 198,4 - 173,1). E o 'x$\in${[}0;5)' torna a região ambígua. Testei outras leituras e nenhuma dá alternativa limpa: {[}0,2{]}?{[}2,4{]} dá 31,3$\pi$ e |x$^{4}$-4\^{}x| em {[}0,5{]} dá 159,7$\pi$.%
\end{itemize}%
\paragraph{Pontos de dúvida}%
\begin{itemize}%
\item%
\textit{(key)} O gabarito preliminar anula a questão sem dizer por quê. Causas prováveis: (i) o valor exato, V $\approx$ 25,28$\pi$, não corresponde a nenhuma alternativa, e só com ln 2 $\cong$ 0,7 se chega a $\approx$27$\pi$, já que o resultado depende muito dessa aproximação; (ii) o enunciado é ambíguo, porque 'x$\in${[}0;5)' junto com 'área delimitada pelas intersecções' não diz com clareza se a região é só {[}2,4{]} ou se inclui os trechos {[}0,2{]} e {[}4,5). Nada disso muda a nota, pois qualquer alternativa é aceita. Vale confirmar no gabarito definitivo se a anulação foi mantida.%
\item%
\textit{(statement)} O intervalo {[}0;5) não influi na leitura natural, que é a região entre as interseções x=2 e x=4. Se a região fosse a de todo o {[}0;5), o volume seria outro: 31,3$\pi$ para {[}0,4{]} ou 159,7$\pi$ para {[}0,5). Nenhuma dessas leituras dá alternativa válida, o que confirma que a leitura {[}2,4{]} da resolução era a pretendida.%
\end{itemize}%
\section*{Questão 13}%
\begin{longtable}{@{}p{0.06\textwidth} p{0.33\textwidth} p{0.33\textwidth} p{0.1\textwidth} p{0.08\textwidth}@{}}\toprule Item & Resolução & Gabarito & Veredito & Certo \\ \midrule
13 & (C) $\alpha$$^2$ (com ? = -$\alpha$ e ? = $\alpha$$^2$ no ponto (2,2)) & C & igual & ambos \\
\bottomrule\end{longtable}%
\begin{itemize}%
\item%
\textbf{13)} O enunciado da transcrição confere com a página 4 do PDF (prova Amarela, Questão 13). Refiz a conta. Com x = 4/y, ? = -4?/y$^2$ e ? = 8?$^2$/y$^3$ - 4ÿ/y$^2$. Como ? = $\alpha$ e ÿ = 0, em x = 2 temos y = 2 e ? = 8$\alpha$$^2$/8 = $\alpha$$^2$. A derivação implícita da resolução (?y + 2?? + xÿ = 0, logo 2? - 2$\alpha$$^2$ = 0) chega ao mesmo valor. Também conferi numericamente, por diferenças finitas, com x(t) = 4/($\alpha$t + 2) para $\alpha$ = 1, -3 e 0,5: deu ?(0) $\approx$ $\alpha$$^2$ nos três casos. A resolução e o gabarito dão a mesma resposta, (C), e ela está correta.%
\end{itemize}%
\paragraph{Pontos de dúvida}%
\begin{itemize}%
\item%
\textit{(statement)} A resolução lê "velocidade de y" como a taxa dy/dt = $\alpha$ constante, o que dá ÿ = 0, e lê "aceleração da abscissa" como d$^2$x/dt$^2$. É a leitura natural, e é com ela que se chega ao gabarito C. Mesmo assim, o enunciado não define essas expressões de forma explícita.%
\item%
\textit{(solution)} O parágrafo de análise dimensional está um pouco confuso: começa falando de $\alpha$$^2$/y e depois passa para 8$\alpha$$^2$/y$^3$. Também não deixa claro que $\alpha$$^2$ só é numericamente igual a 8$\alpha$$^2$/y$^3$ porque y = 2. Isso não afeta o resultado, mas o texto poderia ser mais limpo. O argumento de paridade (inverter o tempo troca $\alpha$ por -$\alpha$ e não muda a aceleração) está correto.%
\end{itemize}%
\section*{Questão 14}%
\begin{longtable}{@{}p{0.06\textwidth} p{0.33\textwidth} p{0.33\textwidth} p{0.1\textwidth} p{0.08\textwidth}@{}}\toprule Item & Resolução & Gabarito & Veredito & Certo \\ \midrule
14 & (B) 128/15 & B & igual & ambos \\
\bottomrule\end{longtable}%
\begin{itemize}%
\item%
\textbf{14)} Conferi a transcrição com o PDF original (prova Amarela, pág. 5/12), e ela está fiel. As curvas se encontram onde x\^{}4 {-} 4x\^{}2 = 0, ou seja, em x = {-}2, 0 e 2. No intervalo {[}{-}2,2{]}, g {-} f = x\^{}2(4 {-} x\^{}2) $\geq$ 0. Em x = 0 a raiz é dupla, então as curvas só se tocam ali e a diferença não troca de sinal. A área é A = $\int$\_\{{-}2\}\^{}\{2\}(4x\^{}2 {-} x\^{}4)dx = 2(32/3 {-} 32/5) = 128/15 $\approx$ 8,533. Refiz a conta com frações exatas e com integração numérica (ponto médio, n = 200000 $\to$ 8,5333333), e os dois resultados batem. A verificação de ordem de grandeza da resolução também está certa: o mínimo de f é {-}9/4, em x$^2$ = 3/2, e o retângulo tem área 4$\cdot$6,25 = 25.%
\end{itemize}%
\paragraph{Pontos de dúvida}%
\begin{itemize}%
\item%
\textit{(statement)} A frase 'área delimitada pelas intersecções das curvas' está redigida de forma imprecisa. A leitura natural, que é a adotada, é a área da região limitada pelas duas curvas entre x = {-}2 e x = 2. Essa região tem dois lóbulos simétricos que só se tocam na origem, onde as curvas são tangentes. Nenhuma outra leitura razoável leva a outra alternativa.%
\end{itemize}%
\section*{Questão 15}%
\begin{longtable}{@{}p{0.06\textwidth} p{0.33\textwidth} p{0.33\textwidth} p{0.1\textwidth} p{0.08\textwidth}@{}}\toprule Item & Resolução & Gabarito & Veredito & Certo \\ \midrule
15 & (B) 3/5 & B & igual & ambos \\
\bottomrule\end{longtable}%
\begin{itemize}%
\item%
\textbf{15)} Conferi o enunciado na página 4 do PDF (impressa como 'Página 5/12', prova Amarela) e a transcrição está fiel, alternativas incluídas. Refiz a conta: P(A)=P(B)=p, P(C)=p/3, P(D)=P(E)=2p/3. A soma dá 11p/3=1, logo p=3/11 e P(D)=2/11. Mantida a proporção, a razão A:D no segundo turno é 3:2, então P2(A)=3/(3+2)=3/5. Isso corresponde à alternativa (B), a mesma do gabarito. A resolução está correta e bem justificada. As alternativas falsas foram identificadas corretamente: 3/11 é a chance de A sem renormalizar, 5/11 é P(A)+P(D) e 2/5 é a chance de D.%
\end{itemize}%
\paragraph{Pontos de dúvida}%
\begin{itemize}%
\item%
\textit{(statement)} Supor que as cinco probabilidades somam 1 equivale a supor que um candidato sempre vence no primeiro turno. Isso não combina com uma eleição em dois turnos, em que pode não haver vencedor no primeiro. Para o resultado, porém, essa suposição não importa: P2(A) depende só da razão P(A):P(D) = p : 2p/3 = 3:2, que vale qualquer que seja a normalização. A resposta 3/5 não muda. A resolução poderia ter registrado que o cálculo de p=3/11 nem é necessário.%
\item%
\textit{(statement)} A expressão 'mantendo{-}se proporcionalmente as chances de ganhar' não é definida com rigor. A leitura natural, que é a adotada, é manter a razão entre as chances de A e D e renormalizar para que somem 1, o que equivale a condicionar ao evento 'vence A ou D'. Não encontrei outra leitura plausível que leve a outra alternativa.%
\end{itemize}%
\section*{Questão 16}%
\begin{longtable}{@{}p{0.06\textwidth} p{0.33\textwidth} p{0.33\textwidth} p{0.1\textwidth} p{0.08\textwidth}@{}}\toprule Item & Resolução & Gabarito & Veredito & Certo \\ \midrule
16 & (D) 81/16 & D & igual & ambos \\
\bottomrule\end{longtable}%
\begin{itemize}%
\item%
\textbf{16)} Confere com o enunciado original (página 4 do PDF, rotulada 'Página 5/12'): a transcrição e as alternativas estão corretas. Refiz a conta: as duas séries geométricas convergem, com soma de 1/3\^{}n igual a (1/3)/(2/3) = 1/2 e soma de 1/2\^{}n igual a (1/2)/(1/2) = 1. Então S = 3/2 e S\^{}4 = 81/16 = 5,0625, que é a alternativa (D). Na verificação da resolução, a soma parcial dá 43/36 $\approx$ 1,19. As alternativas valem aproximadamente: (A) 1,186; (B) 0,886; (C) 4,9375; (E) 1,311. A diferença 81/16 - 79/16 = 2/16 também está certa. Resolução e gabarito estão de acordo e corretos.%
\end{itemize}%
\section*{Questão 17}%
\begin{longtable}{@{}p{0.06\textwidth} p{0.33\textwidth} p{0.33\textwidth} p{0.1\textwidth} p{0.08\textwidth}@{}}\toprule Item & Resolução & Gabarito & Veredito & Certo \\ \midrule
17 & (A) ({-}1)\^{}n$\cdot$(x{-}1)(x{-}2)$\cdot$$\cdot$$\cdot$(x{-}n)/n! & A & igual & ambos \\
\bottomrule\end{longtable}%
\begin{itemize}%
\item%
\textbf{17)} A resolução e o gabarito dão a mesma resposta. Conferi o enunciado no PDF original (pág. 4 do arquivo, impressa como 5/12, prova Amarela): a transcrição está fiel. Pela relação de Stifel para x real, C(x,k) = C(x{-}1,k) + C(x{-}1,k{-}1). Com ela, a soma $\Sigma$\_\{k=0\}\^{}n ({-}1)\^{}k C(x,k) se cancela em cadeia e sobra ({-}1)\^{}n C(x{-}1,n) = ({-}1)\^{}n (x{-}1)$\cdot$$\cdot$$\cdot$(x{-}n)/n!, que é a alternativa (A). Conferi também com frações exatas, para n de 1 a 9 e x $\in$ \{{-}3, 0, 1/2, 2, 7/3, 5, 11\}: a igualdade vale em todos os casos. As verificações da resolução (n=1, n=2, x=0) e a eliminação das outras alternativas por grau e por contraexemplo estão corretas.%
\end{itemize}%
\paragraph{Pontos de dúvida}%
\begin{itemize}%
\item%
\textit{(statement)} O enunciado chama de "Série" e "polinômio" uma soma finita de n+1 termos. Não fica ambíguo, porque a expressão é claramente finita, mas a palavra "Série" pode confundir.%
\item%
\textit{(statement)} A alternativa (B), com numerador (x{-}1)(x{-}2)$\cdot$$\cdot$$\cdot$(x+n), mistura fatores decrescentes e crescentes, então a notação não está bem definida. Isso não afeta a resposta, porque (A) é a única correta sem ambiguidade.%
\end{itemize}%
\section*{Questão 18}%
\begin{longtable}{@{}p{0.06\textwidth} p{0.33\textwidth} p{0.33\textwidth} p{0.1\textwidth} p{0.08\textwidth}@{}}\toprule Item & Resolução & Gabarito & Veredito & Certo \\ \midrule
Questão 18 & (A) a = 1/31 & A & igual & ambos \\
\bottomrule\end{longtable}%
\begin{itemize}%
\item%
\textbf{Questão 18)} Conferi contra o PDF original (prova Amarela, página 6/12) e a transcrição está correta. Nos pontos de interseção vale (a+1)x$^2$ = 8. A condição de ortogonalidade f'(x)$\cdot$g'(x) = (2ax)(-2x) = -4ax$^2$ = -1 leva a 32a = a + 1, logo a = 1/31. Na verificação, com x$^2$ = 31/4, f = g = -27/4 e o produto das inclinações dá -1. Refiz também as contas das alternativas descartadas: B dá -864/31, C dá -136/11, D dá -16/9 e E dá -32/3, todas corretas. A resolução e o gabarito concordam e estão certos.%
\end{itemize}%
\paragraph{Pontos de dúvida}%
\begin{itemize}%
\item%
\textit{(statement)} A alternativa (E) vem escrita como 7/14, sem simplificar (= 1/2). Isso não afeta a resposta e é só uma curiosidade da redação do enunciado.%
\end{itemize}%
\section*{Questão 19}%
\begin{longtable}{@{}p{0.06\textwidth} p{0.33\textwidth} p{0.33\textwidth} p{0.1\textwidth} p{0.08\textwidth}@{}}\toprule Item & Resolução & Gabarito & Veredito & Certo \\ \midrule
19 & (A) -ln($\surd$2/2) + $\pi$/4 - 3$\pi$$^2$/32 $\approx$ 0,2067 & A & igual & ambos \\
\bottomrule\end{longtable}%
\begin{itemize}%
\item%
\textbf{19)} Conferi o enunciado no PDF original (página 6/12, prova Amarela). A transcrição está fiel, inclusive as cinco alternativas. A resolução está correta. Ela usa cot$^2$x = csc$^2$x - 1 e integra por partes, chegando à primitiva F(x) = -x$\cdot$cot x + ln(sen x) - x$^2$/2. Derivando F, volta{-}se ao integrando. Calculando F($\pi$/2) - F($\pi$/4), sai -ln($\surd$2/2) + $\pi$/4 - 3$\pi$$^2$/32. Fiz a integral numericamente pela regra de Simpson com n = 20000 e obtive 0,2066963411, igual ao valor de (A) até a 16ª casa. As outras alternativas valem: (B) -1,796, (C) 4,076, (D) 3,523 e (E) 0,824. Nenhuma delas é compatível com o valor da integral. A cota superior $\pi$$^2$/16 $\approx$ 0,62 usada para eliminar alternativas também vale. O integrando é decrescente no intervalo, porque g'(x) = cot x$\cdot$(sen x cos x - 2x)/sen$^2$x < 0. Logo seu máximo é $\pi$/4, atingido em x = $\pi$/4.%
\end{itemize}%
\paragraph{Pontos de dúvida}%
\begin{itemize}%
\item%
\textit{(solution)} A resolução diz que o integrando decresce de $\pi$/4 até 0 e usa isso para limitar a integral por $\pi$$^2$/16, mas não prova o decrescimento. A afirmação é verdadeira: g'(x) = cot x$\cdot$(sen 2x/2 - 2x)/sen$^2$x < 0 no intervalo. Além disso, essa cota só serve para eliminar alternativas e não afeta o resultado.%
\end{itemize}%
\section*{Questão 20}%
\begin{longtable}{@{}p{0.06\textwidth} p{0.33\textwidth} p{0.33\textwidth} p{0.1\textwidth} p{0.08\textwidth}@{}}\toprule Item & Resolução & Gabarito & Veredito & Certo \\ \midrule
Questão 20 & (D) R\$ 132.300,00 (150.000 $\times$ 0,9 $\times$ 0,7 = 94.500; 94.500 $\times$ 1,4 = 132.300) & D & igual & ambos \\
\bottomrule\end{longtable}%
\begin{itemize}%
\item%
\textbf{Questão 20)} Conferi o enunciado na página 5 do PDF (prova Amarela, Questão 20), e a transcrição está fiel. Refiz a conta: 150000 $\times$ 0,9 = 135000; 135000 $\times$ 0,7 = 94500; 94500 $\times$ 1,4 = 132300. O fator total é 0,882, e o resultado é a alternativa (D), a mesma do gabarito. A resolução trata as perdas sucessivas de forma composta, que é o correto, e mostra que somar as perdas daria 126.000, valor que não está entre as opções.%
\end{itemize}%
\paragraph{Pontos de dúvida}%
\begin{itemize}%
\item%
\textit{(statement)} O enunciado é ambíguo. Diz que o veículo "custa atualmente R\$ 150.000,00" e depois aplica o acréscimo de 40\% sobre o "valor atual". Lido ao pé da letra, o valor atual seria 150.000, e o preço de venda ficaria em 150.000 $\times$ 1,4 = 210.000, que não está entre as alternativas. Outra leitura seria tomar 150.000 como o valor depois das perdas e voltar ao preço de compra (150.000/0,63 $\approx$ 238.095), o que também não bate com nenhuma opção. A única leitura que leva a uma alternativa trata 150.000 como o valor de aquisição, sobre o qual incidem as perdas. É a leitura que a resolução usa e declara explicitamente, e ela é a mesma do gabarito. Há imprecisão de redação, mas a resposta não muda.%
\end{itemize}%
\section*{Questão 21}%
\begin{longtable}{@{}p{0.06\textwidth} p{0.33\textwidth} p{0.33\textwidth} p{0.1\textwidth} p{0.08\textwidth}@{}}\toprule Item & Resolução & Gabarito & Veredito & Certo \\ \midrule
21 & (E): f(x)=2x$^2$-1, então g(x)=1/(2x$^2$-1), com assíntotas verticais em x=$\pm$$\surd$2/2, ramo central negativo em ? com máximo g(0)=-1 e ramos laterais positivos que descem de +$\infty$ até g($\pm$1)=1 & E & igual & ambos \\
\bottomrule\end{longtable}%
\begin{itemize}%
\item%
\textbf{21)} Refiz a conta. Para x em {[}-1,1{]}, cos(2$\cdot$arccos x) = 2x$^2$-1 (conferi numericamente em vários pontos com python). Logo g = 1/(2x$^2$-1) tem assíntotas em x=$\pm$$\surd$2/2$\approx$$\pm$0,707. Para |x|<$\surd$2/2, g $\leq$ -1, com máximo -1 em x=0, e g vai a -$\infty$ perto das assíntotas. Para $\surd$2/2<|x|$\leq$1, g $\geq$ 1: vai a +$\infty$ perto das assíntotas e vale 1 em x=$\pm$1. Olhando a figura original (p. 5), só o gráfico (E) tem ramo central abaixo do eixo, com o -1 marcado no eixo y, e ramos laterais positivos que terminam na linha tracejada horizontal. (A) e (B) não têm parte negativa, (C) tem os sinais trocados e (D) tem o ramo central positivo e passa de x=1. A resolução descarta cada alternativa corretamente e chega à mesma resposta do gabarito.%
\end{itemize}%
\paragraph{Pontos de dúvida}%
\begin{itemize}%
\item%
\textit{(statement)} As figuras são esboços fora de escala. Em (E), as assíntotas e os ramos laterais estão bem estreitos e perto das marcas $\pm$1, e não se vê com clareza onde fica x=$\pm$$\surd$2/2 em relação a essas marcas. Mesmo assim, a escolha não fica ambígua: (E) é a única alternativa com ramo central negativo, então isso não muda a resposta.%
\item%
\textit{(solution)} A resolução diz que (A) não tem assíntotas verticais. Na figura original, os picos de (A) sobem até perto de retas verticais tracejadas, então essa afirmação vale como leitura do esboço. O argumento que de fato descarta (A) é a falta de parte negativa (g(0)=-1), e esse está certo.%
\end{itemize}%
\section*{Questão 22}%
\begin{longtable}{@{}p{0.06\textwidth} p{0.33\textwidth} p{0.33\textwidth} p{0.1\textwidth} p{0.08\textwidth}@{}}\toprule Item & Resolução & Gabarito & Veredito & Certo \\ \midrule
22 & (D) (1+$\surd$2)/8 $\approx$ 0,302 & D & igual & ambos \\
\bottomrule\end{longtable}%
\begin{itemize}%
\item%
\textbf{22)} Refiz a conta com coordenadas centradas na origem, com b = 4+4$\surd$2. Conferi que os oito lados medem 8. A reta BE é x+y = -4$\surd$2 e a reta CH é y = x+4$\surd$2, então I = BE?CH = (-4$\surd$2, 0). Pela fórmula do determinante (shoelace), a área de IEG dá b$^2$ = 16(1+$\surd$2)$^2$ $\approx$ 93,255 e a área do octógono dá 128(1+$\surd$2) $\approx$ 309,019. A razão é (1+$\surd$2)/8 $\approx$ 0,3018. As outras alternativas valem 0,604, 0,402, 0,201 e 0,172, nenhuma perto desse valor. A resolução e o gabarito (D) estão corretos. O enunciado e a figura transcritos conferem com o original (pág. 7/12 da prova Amarela).%
\end{itemize}%
\paragraph{Pontos de dúvida}%
\begin{itemize}%
\item%
\textit{(statement)} A figura não está em escala: o I desenhado fica bem perto de B e acima do centro, mas o I verdadeiro é (-4$\surd$2, 0), na altura do centro. A resolução usa só o dado escrito (a reta IH passa por C) e não a posição desenhada, que é o certo a fazer. Mesmo assim, ao dizer que a razão é 'coerente com a figura', ela se apoia num desenho que não está em escala.%
\item%
\textit{(solution)} A resolução lê 'a reta que contém IH intersecta o ponto C' como 'C, I e H são colineares', e com isso I = BE?CH. É a leitura natural, mas é uma hipótese de interpretação que não foi declarada.%
\end{itemize}%
\end{document}